\documentclass[12pt]{article}

% \usepackage{fullpage}
\usepackage[margin=2cm]{geometry} % поля 2 см со всех сторон
\usepackage{multicol,multirow}
\usepackage{tabularx}
\usepackage{ulem}
\usepackage[utf8]{inputenc}
\usepackage[russian]{babel}

% Оригиналный шаблон: http://k806.ru/dalabs/da-report-template-2012.tex

\begin{document}

\section*{Лабораторная работа №,1 по курсу дискретного анализа: сортировка за линейное время}

Выполнил студент группы 08-215 МАИ \textit{Матвеев Семен}.

\subsection*{Условие}
\hspace{2ex}
Требуется разработать программу, осуществляющую ввод пар «ключ-значение», их упорядочивание по возрастанию ключа указанным алгоритмом сортировки за линейное время и вывод отсортированной последовательности. \\

Вариант задания: 
\begin{itemize}
    \item Алгоритм: Карманная сортировка
    \item Тип ключа: число от 0 до $2^{64}-1$
    \item Тип значения: строки переменной длины (до 2048 символов)
\end{itemize}

\subsection*{Метод решения}
\hspace{2ex}
Для сортировки набора пар (ключ, значение) применяется карманная сортировка. \\

\textbf{Основная идея}:
\begin{enumerate}
\item Исходный массив разбивается на несколько корзин (buckets) на основе значения ключа.
\item Количество корзин выбирается равным $\lfloor\sqrt{n}\rfloor$, где $n$ — число элементов. Такой выбор даёт компромисс между количеством корзин и размером каждой из них.
\item Для каждого элемента вычисляется номер корзины по формуле:
\[
\text{bucketIndex} = \left\lfloor \frac{\text{key} - \min}{1 + \frac{\max - \min}{\text{buckets}}} \right\rfloor
\]
Формула обеспечивает равномерное распределение элементов по корзинам при равномерном распределении ключей (оно является требованием к данным для эффективной работы bucket sort).
\item Каждая корзина сортируется независимо. Для этого используется сортировка слиянием (merge sort), гарантирующая устойчивость и сложность $O(k \log k)$ в худшем случае для корзины размера $k$.
\item После сортировки всех корзин элементы собираются обратно в единый массив в порядке возрастания ключей.
\end{enumerate}
При равномерном распределении ключей размер каждой корзины составляет примерно $\sqrt{n}$, поэтому общая трудоёмкость близка к $O(n)$, так как сортировка каждой корзины даёт $O(\sqrt{n} \log \sqrt{n})$, а сумма по всем корзинам — $O(n \log \sqrt{n}).

\subsection*{Описание программы}
\hspace{2ex}
\textbf{Шаблон класса \text{TVector<T>}} — динамический массив с ручным управлением памятью. Поддерживает копирование, перемещение, доступ по индексу, изменение размера, добавление элемента в конец. Особенность: в методе \text{push\_back} используется \text{placement new}, что заметно уменьшило время выполнения программы в процессе отладки. \\

\textbf{Перегруженные операторы ввода и вывода} — для \text{TVector<T>} определён вывод всех элементов подряд. Для \text{TVector<char>} есть специализированный оператор ввода, который читает последовательность непробельных символов. \\

\textbf{Структура \text{TItem}} — хранит ключ (\text{size\_t key}) и значение (\text{TVector<char> value}), имеет перегруженные операторы сравнения. \\

\textbf{Вспомогательные функции} — шаблонные \text{Min} и \text{Max}, \text{GenerateRandomInput} (для генерации случайных входных данных), шаблонная \text{RunBenchmark} (для теста производительности на случайных выходных данных), BubbleSort (обменная сортировка для сравнения с \text{BuckerSort} в рамках \text{RunBenchmark}). \\

\textbf{Функции сортировки} — \text{Merge} (слияние двух отсортированных половин), \text{MergeSort} (итеративная сортировка слиянием всего массива), \text{BucketSort} (основная функция карманной сортировки). \\

\textbf{Функция \text{main}} — читает входные данные до конца потока, сортирует их вызовом \text{BucketSort}, затем выводит каждую пару на экран. Опционально, раскомментировав блоки кода, можно использовать случайный ввод вместо ручного пользовательского и провести бенчмарк функции сортировки. \\

\subsection*{Дневник отладки}

\begin{enumerate}
\item Реализовал каркас программы - ввод данных в виде набора пар (ключ, значение), сохранение их в \text{std::vector}, их сортировка с помощью \text{std::sort}, вывод на экран. Пришлось осознать, что в выводе надо использовать символы табуляции.

\item Реализовал функцию сортировки \text{MergeSort} и вспомогательные функции \text{Merge}, \text{Min} и \text{Max}. Возникали ошибки выхода за границы массива из-за неверных индексов в циклах. Исправлено изменением логики использования переменных \text{left}, \text{mid}, \text{right}.

\item Реализовал функцию \text{BucketSort}. Возникали ошибки  выхода за границы массива из-за некорректного вычисления переменной \text{bucketIndex}. Я ориентировался на нерекурсивную реализацию карманной сортировки в конспектах ИТМО, но там эта формула была чрезвычайно странной. Сначала боролся с переполнением, потом с логикой вычисления (классическая нормализация и добавление 1, чтобы гарантировать попадание максимального ключа в последнюю корзину).

\item Реализовал вспомогательные функции \text{GenerateRandomInput} (которая также помогала при отладке), \text{RunBenchmark} и \text{{BubbleSort} для проведения тестов производительности функции \text{BucketSort}.

\item Самое сладкое. Было необходимо избавиться от использования \text{STL} (в моем случае \text{std::pair, std::vector, std::string}). Для этого были реализованы классы \text{TItem}, \text{TVector} и все их методы. Также понадобилось перегрузить \text{operator>>} и \text{operator<<} для ввода и вывода объектов класса \text{TVector} соответственно. Ошибок было несчетное множество, все уже и не упомнишь. Могу лишь сказать, что я в какой-то момент напрягся настолько, что единолично увеличил уровень общемировой напряженности, чем спровоцировал очередной конфликт на Ближнем Востоке. За что искренне прощу прощения.

\end{enumerate}

\subsection*{Тест производительности}
\hspace{2ex}
Для оценки эффективности реализованных алгоритмов был проведён эксперимент на случайных данных. Генерация входных наборов выполнялась функцией \text{GenerateRandomInput} с фиксированной строкой-значением \text{'stroka'} и равномерно распределёнными 64-битными ключами (диапазон $[0, 2^{64}-1]$). \\

Время сортировки измерялось с помощью \text{std::chrono::high\_resolution\_clock}; для каждого объёма данных проводилось три запуска, результат усреднялся. \\

Исследовались следующие алгоритмы:
\begin{itemize}
\item \text{BucketSort} — карманная сортировка с количеством корзин $\lfloor\sqrt{n}\rfloor$ и внутренней сортировкой слиянием;
\item \text{BubbleSort} — пузырьковая сортировка (включена в программу для демонстрации контраста сложности).
\end{itemize}
Из-за квадратичной сложности пузырьковой сортировки её тестирование выполнялось только для небольших $n$ (до 50,000), чтобы время измерений оставалось приемлемым. \\

Результаты измерений приведены в таблице.

\begin{center}
\begin{tabular}{|c|c|c|}
\hline
\textbf{Количество элементов} & \textbf{BucketSort, мс} & \textbf{BubbleSort, мс} \\
1,000 & 0.69 & 8.23 \\
5,000 & 3.90 & 207.26 \\
10,000 & 8.33 & 841.87 \\
50,000 & 41.64 & 21038.40 \\
100,000 & 89.26 & – \\
500,000 & 500.58 & – \\
1,000,000 & 1077.96 & – \\
\hline
\end{tabular}
\end{center}

Для \text{BucketSort} прослеживается близкая к линейной зависимость времени от размера входных данных: при увеличении $n$ в 10 раз (с 100,000 до 1,000,000) время выполнения увеличивается в 12.07 раза, что согласуется с теоретической оценкой $O(n)$ при равномерном распределении ключей.  \\

Пузырьковая сортировка демонстрирует ожидаемый квадратичный рост: при увеличении $n$ в 5 раз (с 1,000 до 5,000) время выполнения увеличивается в 25 раз, что соответствует $O(n^2)$. \\

Полученные результаты подтверждают корректность реализации и соответствие заявленной сложности.

\subsection*{Недочёты}

\begin{itemize}
\item При вычислении индекса корзины используется деление с плавающей точкой (\text{long double}), что может привести к небольшим погрешностям, с которыми, вроде бы, получилось совладать.
\item Не обрабатывается случай пустого входного потока — программа завершится без вывода, что допустимо, но можно было бы добавить проверку.
\item Отсутствует явная проверка на превышение максимальной длины строки (2048 символов), так как автор программы доверяет условию. В текущей реализации строка может быть любой длины, что потенциально может привести к исчерпанию памяти.
\end{itemize}

\subsection*{Выводы}
\hspace{2ex}
В результате выполнения лабораторной работы реализована программа сортировки пар «ключ-значение» методом карманной сортировки. Ключи — 64-битные целые числа, значения — строки переменной длины. Все компоненты программы разработаны без использования STL, что позволило глубже изучить управление памятью и алгоритмы. \\

В процессе отладки были устранены ошибки выхода за границы массивов, некорректного вычисления индексов корзин и ввода строк. Применение \text{placement new} позволило сократить накладные расходы на копирование. \\

Тестирование на равномерно распределённых данных подтвердило близкую к линейной зависимость времени работы от объёма входных данных ($O(n)$), что соответствует теоретической оценке для карманной сортировки. Для сравнения пузырьковая сортировка показала ожидаемый квадратичный рост, что косвенно подтверждает корректность измерений. \\

Таким образом, задача решена в полном объёме, полученные навыки могут быть использованы при разработке эффективных алгоритмов сортировки в условиях ограниченных ресурсов. \\

\end{document}